% =====================================================================
%  1 INTRODUÇÃO — "Qual é o problema da pesquisa?"
%  Ordem: contextualização -> problema -> questão de pesquisa (em forma de
%  pergunta) -> objetivo geral e específicos -> justificativa (pode virar
%  seção própria, como no modelo do INCRA) -> estrutura do trabalho.
%  Escrita SEMPRE impessoal (Dica 1): "propõe-se", "foi realizado".
% =====================================================================
\section{Introdução}

\escrever{contextualização: luvas sensorizadas capturam o movimento das mãos, mas raramente medem forças de contato nas pontas dos dedos; sensores resistivos de baixo custo sofrem com histerese, deriva e baixa repetibilidade. Citar a literatura, p.\,ex. \citeonline{dahiya2010tactile} e \citeonline{vanwegen2023haptics}.}

\escrever{estado da técnica em 1 parágrafo: sensores táteis magnéticos e o eFlesh \cite{pattabiraman2025eflesh}, com os números de desempenho reportados.}

\escrever{problema: o eFlesh foi demonstrado em blocos de 40$\times$40~mm, com cinco magnetômetros; não há relato da versão em escala de ponta de dedo humano nem da integração de múltiplas unidades em uma luva.}

\escrever{questão de pesquisa em forma de pergunta. Exemplo de estrutura: É possível miniaturizar ... mantendo ...?}

\escrever{objetivo geral (o MESMO texto que irá para o Resumo — Dica 7).}

Como objetivos específicos, propõem-se:
\begin{enumerate}[label=\alph*)]
  \item \escrever{projetar e fabricar por manufatura aditiva a microestrutura em escala de ponta de dedo;}
  \item \escrever{caracterizar em bancada a resposta a força normal, cisalhante e posição de contato;}
  \item \escrever{calibrar e comparar modelos de regressão e \ac{MLP} com validação temporal;}
  \item \escrever{embarcar o modelo no microcontrolador e fechar a malha com feedback vibrotátil;}
  \item \escrever{validar o conjunto em tarefas com objetos reais.}
\end{enumerate}

\escrever{justificativa: relevância técnica e social (caso de uso que justifica o feedback no mundo real).}

\escrever{estrutura do trabalho: Este trabalho está organizado da seguinte forma: a Seção~\ref{sec:estado-arte} apresenta ...; a Seção~\ref{sec:referencial} ...; a Seção~\ref{sec:metodo} ...; a Seção~\ref{sec:resultados} ...; e a Seção~\ref{sec:conclusoes} ...}
